\section*{الفصل \ref{ch:g9:ions} --- الأيونات والمحاليل الأيونية}

\begin{solution}{exo:g9:ions:1}
\omterm{def:g7:atoms-and-molecules:atom}{ذرة} أو مجموعة من \omterm{def:g7:atoms-and-molecules:atom}{الذرات} فقدت \omterm{def:g9:inside-the-atom:nucleus}{إلكترونات} أو اكتسبتها فصارت تحمل شحنة.
\omterm{def:g9:ions:ion}{والكاتيون} موجب (فقد \omterm{def:g9:inside-the-atom:nucleus}{إلكترونات})، \omterm{def:g9:ions:ion}{والأنيون} سالب
(اكتسب \omterm{def:g9:inside-the-atom:nucleus}{إلكترونات}).
\end{solution}

\begin{solution}{exo:g9:ions:2}
\ce{Mg^{2+}}: 12 بروتونًا، و10 \omterm{def:g9:inside-the-atom:nucleus}{إلكترونات}.
\end{solution}

\begin{solution}{exo:g9:ions:3}
\ce{KCl}، \ce{MgCl2}، \ce{CaO}.
\end{solution}

\begin{solution}{exo:g9:ions:4}
الماء المالح يحوي \omterm{def:g9:ions:ion}{أيونات}، \ce{Na+(aq)} و\ce{Cl-(aq)}، تتحرك وتحمل
الشحنة. أما \omterm{def:g7:atoms-and-molecules:molecule}{جزيئات} السكر فلا تحمل شحنة.
\end{solution}

\begin{solution}{exo:g9:ions:5}
\ce{CuSO4}، \ce{AlCl3}، \ce{Na2CO3}.
\end{solution}

\begin{solution}{exo:g9:ions:6}
الحديد (III)، \ce{Fe^{3+}}: \ce{Fe^{3+}(aq) + 3OH-(aq) -> Fe(OH)3(s)}.
\end{solution}

\begin{solution}{exo:g9:ions:7}
الكبريتات: إلكترونان زائدان. النترات: \omterm{def:g9:inside-the-atom:nucleus}{إلكترون} واحد.
\end{solution}

\begin{solution}{exo:g9:ions:8}
الأبيض مع نترات الفضة: \ce{Cl-}. والأزرق مع هيدروكسيد الصوديوم:
\ce{Cu^{2+}}. إنه كلوريد النحاس (II)، \ce{CuCl2}.
\end{solution}

\begin{solution}{exo:g9:ions:9}
أربعة: واحد من كل جهة (يسارًا، ويمينًا، وأعلى، وأسفل).
\end{solution}

\begin{solution}{exo:g9:ions:10}
\ce{Fe^{2+}(aq) + 2OH-(aq) -> Fe(OH)2(s)}.
\end{solution}

\begin{solution}{exo:g9:ions:11}
كلوريد الزنك: \omterm{def:g9:ions:ion}{فأيونات} الصوديوم لا تعطي راسبًا مع هيدروكسيد الصوديوم، أما
\omterm{def:g9:ions:ion}{أيونات} الزنك فتعطي راسبًا أبيض:
\ce{Zn^{2+}(aq) + 2OH-(aq) -> Zn(OH)2(s)}.
\end{solution}

\begin{solution}{exo:g9:ions:12}
\ce{CaCl2}: 2000 \omterm{def:g9:ions:ion}{أيون} كلوريد. والشحنات: $1000 \times (+2) + 2000 \times
(-1) = 0$.
\end{solution}

\begin{solution}{pb:g9:ions:1}
\textbf{1.} الحديد (II)، \ce{Fe^{2+}}: \ce{Fe(OH)2}.

\textbf{2.} الحديد (III)، \ce{Fe^{3+}}: \ce{Fe(OH)3}.

\textbf{3.} في \ce{Fe^{2+}} \omterm{def:g9:inside-the-atom:nucleus}{إلكترون} زائد: $26 - 2 = 24$ إلكترونًا، مقابل
$26 - 3 = 23$ في \ce{Fe^{3+}}.

\textbf{4.} \ce{Fe^{3+}(aq) + 3OH-(aq) -> Fe(OH)3(s)}.

\textbf{5.} يخرج حاملًا \omterm{def:g9:ions:ion}{أيونات} الحديد (II)؛ وبملامسة \omterm{def:g4:air-a-mixture-of-gases:air}{الهواء} تصير \omterm{def:g9:ions:ion}{أيونات}
حديد (III)، تكوّن الجسم الصلب البرتقالي.

\textbf{6.} \ce{Fe(OH)3}: $(+3) + 3 \times (-1) = 0$.

\textbf{7.} ليس لحظة خروجه: \omterm{def:g9:ions:ion}{فأيونات} الحديد ذائبة وتمر عبر المرشِّح. أما بعد
تكوّن الجسم الصلب البرتقالي فنعم: فهو \omterm{def:g9:ions:precipitate}{راسب} يحجزه
المرشِّح.

\textbf{8.} لا: إنه \omterm{def:g6:pure-substances-and-mixtures:mixture}{خليط} (ماء \omterm{def:g9:ions:ion}{وأيونات} ذائبة). وهو مسحوبًا حديثًا ورائقًا
\omterm{def:g6:pure-substances-and-mixtures:homogeneous}{خليط متجانس}.

\textbf{9.} $56 \times 1.67 \times 10^{-27} = \qty{9.35e-26}{kg}$.

\textbf{10.} $\qty{0.30}{mg} = \qty{0.30e-6}{kg} = \qty{3.0e-7}{kg}$.

\textbf{11.} يختلف \omterm{def:g9:ions:ion}{الأيون} عن \omterm{def:g7:atoms-and-molecules:atom}{الذرة} بإلكترونين أو ثلاثة، كتلتها مهملة إلى
جانب \omterm{def:g9:inside-the-atom:nucleus}{النواة}.

\textbf{12.} $3.0 \times 10^{-7} / 9.35 \times 10^{-26} \approx 3.2
\times 10^{18}$ \omterm{def:g9:ions:ion}{أيون} حديد في اللتر.
\end{solution}
